\subsection{Лекция 19 (14.03.26) Задача о естественной конвекции. Метод дополнений Шура. МКЭ}
В тесте \cvar{[heat-fem-simple]} из файла \ename{heat_fem_simple_test.cpp}
реализована конечноэлементная схема решения нестационарной задачи о естественной
конвекции в закрытом объёмы с цилиндрическим нагревателем.
Используются MINI-элементы:
давление аппроксимируется на линейных треугольниках,
скорости -- на треугольниках с bubble-точкой.
Решения СЛАУ осуществляется методом простой
итерации с предобуславливателем Шура.

Безразмерная определяющая система уравнений приведена ниже:

\begin{equation}
\label{eq:natural_convection_y}
\begin{aligned}
& \dfr{u}{t} + \vec u\cdot\nabla u = \frac{1}{\Ren} \nabla^2 u - \dfr{p}{x}, \\
& \dfr{v}{t} + \vec u\cdot\nabla v = \frac{1}{\Ren} \nabla^2 v - \dfr{p}{y} + T, \\
& \nabla \cdot \vec u = 0, \\
& \dfr{T}{t} + \vec u\cdot\nabla T = \frac{1}{\Pen} \nabla^2 T
\end{aligned}
\end{equation}
Характерная скорость определена как
$u_0 = \sqrt{gL(T_1 - T_0) \beta}$, где $\beta$ -- коэффициент сжимаемости, $g$ -- ускорение свободного падения, $L$ -- характерный масштаб длины, $T_0, T_1$ -- максимальная и минимальная температуры.

Необходимо:
\begin{enumerate}
\item  Провести расчёт до установления. Нарисовать динамику полученных полей скорости и температур.
\item  Подсчитать коэффициент теплоотдачи: локальный  (на разные моменты времени)
       $$\rm Nu(\alpha) = \frac{1}{\rm Pe} \dfr{T}{n} (\alpha)$$,
       зависящий от угловой координаты $\alpha$ точки на нагревателе,
       и интегральный
       $$\overline{\rm Nu}(t) = \frac{1}{\rm Pe} \int_{\gamma}\dfr{T}{n}(t)\, ds$$.
\item  Оптимизировать процедуры \cvar{compute_residual}, \cvar{use_preconditioner}
       используя операции с блоками и отказавшись от сборки глобальных матрицы. Замерить время работы оптимизированных процедур
\item  В методе \cvar{step} вместо простого итерационного процесса \secref{sec:it_def_corr} использовать BiCGStab \secref{sec:bicgstab}.
\end{enumerate}

\paragraph{Рекомендации к программированию}
\begin{itemize}
\item п. 2
Для подсчёта нормальной производной по температуре нобходимо
пройти циклом по всем ячейкам, лежащим около границы нагревателя.
Далее найти температуру в центре этой ячейки $T_c$ (как среднее от температур в узлах этой ячейки).
Мы знаем, что на границе нагревателя $T(r=0.5) = 1.0$.
Отсюда нормальная производная на текущей грани выпишется как
\begin{equation*}
\left.\dfr{T}{n}\right|_\gamma = \frac{1.0 - T_c}{h_c - 0.5},
\end{equation*}
Здесь $h_c$ -- расстояние от центра ячейки до центра нагревателя в точке $(0, 0)$.
Для нахождения интегрального коэффициента теплоотдачи это слагаемое нужно умножить
на площадь грани и добавить в сумму. Для локального -- найти угловую координату по центральной точке грани и приписать значение ей.
Некоторые методы, необходимые для вычисления:
\begin{itemize}
\item \cvar{grid_.boundary_faces()} -- список индексов всех граничных граней, 
\item \cvar{grid_.face_сenter(iface)} -- координтата центра грани,
\item \cvar{grid_.tab_face_cell(iface)} -- две ячейки, инцидентные грани. Если грань граничная, одна из них равна \cvar{INVALID_INDEX}
\item \cvar{grid_.сell_сenter(icell)} -- координтата центра ячейки,
\item \cvar{grid_.tab_cell_point(icell)} -- список узлов, инцидентных ячеке. Необходим для вычисления $T_c$ из массива \cvar{temperature_}.
\end{itemize}

\item п. 3
В методе \cvar{compute_resudual} вычисляется выражение
\begin{equation*}
r = rhs - 
\left(
\begin{matrix}
A & B^T\\
B & Z  \\
\end{matrix}
\right)
\left(
\begin{matrix}
u \\
p
\end{matrix}
\right)
\end{equation*}
Где $Z$ -- это нулевая матрица с модификациями, связанными с граничными условиями первого рода на давление (в одной точке).
В неоптимально написанном методе сначала из блоков собирается глобальная матрица \cvar{mat},
а потом проводится умножение.
Можно обойтись без сборки, если написать
\begin{align*}
&r_u = rhs_u - \mat A u  - \mat B^T p, \\
&r_p = rhs_p - \mat B u  - \mat Z p.
\end{align*}
Учесть, что длины векторов равны:
\begin{cppcode}
// вектора rhs_u, r_u, u, y_u
// Две компоненты скорости, каждая из которых задана
// на трёх узлах треугольников и bubble-точке
size_t nu = 2 * grid_.n_cells() + 2 * grid_.n_points();
// вектора rhs_p, r_p, p, y_p
// Заданы только в узлах треугольных ячеек
size_t np = grid_.n_points();
\end{cppcode}
Для получения подмассива из сплошного вектора можно воспользоваться арифметикой указателей
\begin{cppcode}
const double* u = uvp_.data() + 0;  // без сдвига от начала uvp_
const double* p = uvp_.data() + nu; // сдвиг на nu элементов

// Умножение матрицы на вектор
std::vector<double> Au = A_->mult_vec_p(u);
\end{cppcode}

\item п. 3
В методе \cvar{use_preconditioner} вычисляется выражение
\begin{equation*}
\left(
\begin{matrix}
A & 0\\
B & -S  \\
\end{matrix}
\right)
\left(
\begin{matrix}
I & H B^T\\
0 & I  \\
\end{matrix}
\right)
x = r
\end{equation*}

Сначала матрица левой части собирается в глобальную,
потом вызывается решатель, для нахождения неизвестного вектора $x$ в левой части.
Полученная матрица имеет проблему седловой точки (нули на диагоналях),
поэтому её нельзя решить общими итерационными методами.
Для её решения исплользуется точный решатель \cvar{UmfpackSolver}.
Предлагается решать эту систему с следующем порядке
\begin{equation*}
\left(
\begin{matrix}
A & 0 \\
B & -S
\end{matrix}
\right)
\left(
\begin{matrix}
y_u\\
y_p
\end{matrix}
\right) = 
\left(
\begin{matrix}
r_u\\
r_p
\end{matrix}
\right)
\end{equation*}
откуда
\begin{align*}
&y_u = A^{-1} r_u,\\
&y_p = S^{-1} (B y_u - r_p).
\end{align*}
Далее
\begin{equation*}
\left(
\begin{matrix}
I & HB^T \\
0 & I
\end{matrix}
\right)
\left(
\begin{matrix}
x_u\\
x_p
\end{matrix}
\right) = 
\left(
\begin{matrix}
y_u\\
y_p
\end{matrix}
\right)
\end{equation*}
откуда
\begin{align*}
&x_p = y_p, \\
&x_u = y_u - H B^T x_p
\end{align*}
Для вычисления $y$ воспользоваться матричным решателем.
Сначала точным (\cvar{UmfpackSolver}),
потом, если всё работает, можно перейти на итерационный
(\cvar{AmgcMatrixSolver}) и сравнить время работы.

\item п. 4
Метод BiCGStab с учётом предобуславливателя приведён в \secref{sec:bicgstab}.
На шагах 3 инициализации, а также 4 и 7 -- основных итераций
требуется обратить предобуславливатель. Это нужно делать
написанным методом \cvar{use_preconditioner}.
Поскольку один и тот же предобуславливатель
обращается три раза, следует разделить инициализацию
матричного решателя (делать один раз) и непосредственно решени (делать три раза)
\begin{cppcode}
// инициализация
UmfpackSolver slv_A;
slv_A.set_matrix(A_);
// решения
std::vector<double> r0_prime;
slv_A.solve(r0, r0_prime);
\end{cppcode}
Для сравнения скорости сходимости с методом простой итерации имеет смысл 
уменьшить $\epsilon$ -- условие выхода из итерационного процесса.

\end{itemize}
